% !PRINTED none

\begin{center}
\gesperrt{Siebenter Abschnitt.}

Gruppen bin\"arer linearer Substitutionen.

\rule{4cm}{0.4pt}

\S.\,48.

\gesperrt{Tern\"are orthogonale Substitutionen.}
\end{center}

Es ist nun unsere Aufgabe, aus der Gesammtheit der linearen
Substitutionen engere Gruppen auszusondern und schliesslich
zu endlichen Gruppen zu gelangen. Solche engere Gruppen, die
immer noch unendlich sein k\"onnen, erh\"alt man, wenn man die
Forderung stellt, dass gegebene homogene Functionen der
Variablen invariant bleiben sollen. Wir wollen aber die Aufgabe
nicht in dieser Allgemeinheit weiter verfolgen, sondern gleich
zur Betrachtung des wichtigsten speciellen Falles \"ubergehen. Wir
wollen uns auf tern\"are Substitutionen beschr\"anken und fordern,
dass eine quadratische Form von nicht verschwindender Determinante invariant bleiben soll. Da, wie wir fr\"uher gesehen
haben (Bd. I, \S.\,57), jede solche quadratische Form durch lineare
Transformation in eine Summe von Quadraten verwandelt werden
kann, so beschr\"anken wir das Problem nicht weiter, wenn wir
f\"ur diese quadratische Form die Summe der Quadrate annehmen.
Solche Substitutionen heissen \gesperrt{orthogonal}.

Die Substitution

\weq{1}{(y_1, y_2, y_3) = A (x_1, x_2, x_3)}

ist also orthogonal, wenn die Substitutionscoëfficienten so bestimmt sind, dass die Identit\"at besteht:

\weq{2}{y_1^2 + y_2^2 + y_3^2 = x_1^2 + x_2^2 + x_3^2.}

Wenn man die Ausdr\"ucke (1) in (2) substituirt und die
Coëfficienten entsprechender Glieder einander gleich setzt, so
%CONT>>
